% ===== 原文 pp.71–73（PDF p.81–83） =====
% [原文 p.71]

\subsection{物质扇区的边界条件}

我们可以对周期标量 $A_4$ 施加 Dirichlet 条件或 Neumann 条件。下面依次讨论这两种可能性。

\subsubsection{Dirichlet 条件}

若 $A_4$ 满足 Dirichlet 条件，则 BRST 不变性要求 1-形式 $\phi$ 满足 Neumann 条件。这意味着 $\phi$ 与边界相切的分量是自由的，且满足 $\partial_3\phi_1=\partial_3\phi_2=0$；而分量 $\phi_3$ 在边界上取固定值 $\phi_3=a$。BRST 不变性还要求下列费米子在边界上为零：$\psi_4+\tilde\psi_4$、$\psi_3-\tilde\psi_3$、$\rho_1$、$\rho_2$。实参数 $a$ 与 $A_4$ 的边界值合起来构成一个取值于 $S^1\times\R$ 的参数。这些参数其实是不相关的（irrelevant）：拓扑关联函数并不依赖于它们。为看出这一点，注意平移 $A_4$ 的边界值可以通过给作用量加上如下形式的边界项来实现：
\[
  \int_{x_3=0}\partial_3A_4\,\dd^2x
  =\int_{x_3=0}\left(\delta_Q\rho_3-(\partial_1\phi_2-\partial_2\phi_1)\right)\dd^2x .
\]
我们看到，在相差一个全导数项的意义下，这个边界项是 BRST-正合的，因此不影响关联函数。对 $\phi_3$ 的边界值可以作类似的论证。

% [Tutor Note: p.71 Dirichlet 条件下 $A_4$ 与 $\phi$ 的配对可用分量计数自查：$A_4$ 取 Dirichlet 时，条件是 $\partial_3\phi_1=\partial_3\phi_2=0$（相切两分量法向导数为零）加上 $\phi_3=a$（法向分量固定），共 3 条；而 $A_4$ 取 Neumann 时（p.73）变成 $\phi_1=\phi_2=0$（相切两分量在边界上为零，两条）加 $\partial_3\phi_3=0$（一条），同样是 3 条。另注意 $\psi_4+\tilde\psi_4$、$\psi_3-\tilde\psi_3$ 中的波浪号：文字层会把 $\tilde\psi$ 抽成 $\psi$，符号必须按 p.81 图像核对。]

% [Tutor Note: p.71 边界项恒等式右端的 $\delta_Q\rho_3$ 说明该边界项在相差全导数下 BRST-正合；这里 $\rho$ 是 3d 中 1-形式费米子，$\rho_3$ 是它在法向的分量。核对方式：对作用量加该边界项等价于把 $A_4$ 的边界值平移一个常数，而拓扑关联函数只依赖 BRST 上同调类，故不依赖于该平移。]

% [原文 p.72]
把理论放在 $W=\Sigma\times I$ 上（$\Sigma$ 是二维流形，$I$ 是带 Dirichlet 边界条件的区间）会给出 $\Sigma$ 上一个新的 2d TFT，我们将在第 7 章称它为\emph{修正 A-模型}（modified A-model）（在那里我们将更详细地研究它的性质）。对于带 $U(1)$ 作用的靶空间 $X$，可以把通常的 2d A-模型扭转，得到以 $X$ 为靶空间的修正 A-模型。在我们的情形中，$X=\C$，带有按相位的旋转作用。通常 A-模型中 $\C$-值玻色场 $Z$ 在扭转后变成一个实 1-形式（其两个分量分别是 $Z$ 的实部与虚部）。这就是约化后存活下来的 1-形式 $\phi|_\Sigma$。

除玻色 1-形式 $\phi$ 之外，修正 A-模型还有一个费米 1-形式 $\psi-\tilde\psi$ 以及一对费米 0-形式 $\eta-\tilde\eta$ 与 $\rho$（后者来自 3d 中 1-形式 $\rho$ 的分量 $\rho_3$）。它们的 BRST 变换为
\[
\begin{aligned}
  \delta_Q\phi &= \ii(\psi-\tilde\psi),\\
  \delta_Q(\psi-\tilde\psi) &= 0,\\
  \delta_Q(\eta-\tilde\eta) &= 2d^{\star}\phi,\\
  \delta_Q\rho &= \star\dd\phi .
\end{aligned}
\]
这里 $\star$ 是 2d Hodge 星算子，且 $d^{\star}=\star\dd\star$。

为理解 3d 中边界线算子的范畴，我们需要描述修正 A-模型的边界条件范畴。这相当直接。一类自然的边界条件由在边界上施加
\[
  (a\phi+b\star\phi)|_{\partial W}=0
\]
得到。特殊情形 $b=0$ 与 $a=0$ 分别对应于 2d 的 Dirichlet 条件与 Neumann 条件。由于理论显然具有把 $\phi$ 变到 $\star\phi$ 的对称性，只需考虑 Neumann 条件 $\star\phi|_{\partial W}=0$。BRST 不变性要求 $\star\psi$ 与 $\rho$ 在这类边界上的限制为零。容易看出，不存在非平凡的 BRST 不变边界可观测量（唯一 BRST 不变的费米子 $\psi$ 是 BRST-正合的），因此没有可能以非平凡的方式耦合边界自由度。这意味着边界条件的范畴与平凡 2d TFT 的情形相同，即有限维向量空间复形的范畴。我们可以把它记作 $D^b(\mathrm{Coh}(\bullet))$。

这里有一个重要的微妙之处，它与标量 $A_4$ 的周期为 1 这一事实有关。在区间上约化时，这意味着存在“缠绕扇区”，其中
\[
  \int \dd x^3\partial_3A_4=n,\qquad n\in\Z .
\]
这种缠绕沿边界的连通分支为常数，并且完全不影响 2d 理论。

% [Tutor Note: p.72 为何边界条件的范畴是 $D^b_{\C^*}(\mathrm{Coh}(\bullet))$ 而不是 $D^b(\mathrm{Coh}(\bullet))$（见 p.73）：$A_4$ 的周期为 1 使 $\int\dd x^3\partial_3A_4\in\Z$，这一整数缠绕数沿边界分支守恒，给出每个分支上的额外整数标号；整数标号对应 $\C^*$ 的权（特征标格为 $\Z$），故得到 $\C^*$-等变版本。核对：若丢掉缠绕扇区，就会在 $t=1$ 的 DD 条件下把单位对象误判为 $D^b(\mathrm{Coh}(\bullet))$ 中的对象，与 p.75 的 B-侧结果不符。]

% [原文 p.73]
我们可以通过给每个边界分支引入一个额外的整数标号来纳入它，该标号充当守恒的边界荷。这在数学上等价于说：边界条件的范畴是单点上方 $\C^*$-等变凝聚层的范畴 $D^b_{\C^*}(\mathrm{Coh}(\bullet))$。该范畴的对象是带 $\C^*$ 作用的有限维向量空间复形，其中复形中的微分与 $\C^*$ 作用交换。态射要求保持 $\C^*$ 作用，即具有零 $\C^*$ 荷。

\subsubsection{Neumann 条件}

若 $A_4$ 满足 Neumann 条件 $\partial_3A_4=0$，则 BRST 不变性要求 $\phi$ 满足 Dirichlet 条件。也就是说，$\phi$ 在边界上的限制必须为零，且 $\phi_3$ 必须满足 $\partial_3\phi_3=0$。这个边界条件没有任何参数。

在区间上约化会给出以 $A_4$ 与 $\phi_3$ 为玻色场的 A-模型。例如，可以通过考察 3d BPS 方程 $\dd A_4+\star\dd\phi=0$、并限制到满足 $\phi_1=\phi_2=0$ 且 $A_4$ 与 $\phi_3$ 不依赖于 $x^3$ 的场位形来看出这一点。对这样的场位形，BPS 方程变成以 $S^1\times\R\simeq\C^*$ 为靶空间的全纯瞬子方程。从辛的观点看，带标准 Kähler 形式的 $\C^*$ 同构于 $T^*S^1$。因此这一情形中边界线算子的范畴是 $T^*S^1$ 的 Fukaya-Floer 范畴。由于该范畴是作为 3d TFT 中边界线算子的范畴出现的，它必定具有幺半结构。尽管此范畴表面上与规范扇区中 Dirichlet 边界条件下的范畴相同，我们将看到其幺半结构完全不同，它由 $T^*S^1$ 的（某个版本的）Fukaya-Floer 范畴与 $S^1$ 的可构造导出范畴之间的等价所诱导 [43]。特别地，单位对象（即不可见的边界线算子）不同，它对应于 $T^*S^1$ 的零截面配以平凡的秩 1 局部系统。这说明：2d TFT 中膜上的幺半结构依赖于该 2d TFT 实现为 3d TFT 在区间上紧化的方式。

% [Tutor Note: p.73 “BPS 方程变成全纯瞬子方程”一步可显式核对：取 $\phi_1=\phi_2=0$、$\partial_3=0$、$A_4=A_4(x_1,x_2)$、$\phi_3=\phi_3(x_1,x_2)$，则 $\dd A_4+\star\dd\phi=0$ 化为 $\partial_1A_4=-\partial_2\phi_3$、$\partial_2A_4=\partial_1\phi_3$，即 $A_4-\ii\phi_3$ 关于 $z=x_1+\ii x_2$ 全纯（按 $\star$ 的取向约定可能整体取复共轭）。这正是以 $S^1\times\R\simeq\C^*$ 为靶空间的全纯瞬子方程，也解释了为何靶空间是 $\C^*$ 而不是 $S^1\times\R$ 的某个扭曲版本。]

\subsection{\texorpdfstring{$t=\ii$ 与 $t=1$}{t=i 与 t=1} 时面算子的电磁对偶}

现在我们可以描述 4d 电磁对偶如何作用在上述各种边界条件上。由于规范扇区与物质扇区各自都可以取 Dirichlet 或 Neumann 条件，共有四种可能性需要考虑。

从 3d 的观点看，4d 电磁对偶相当于把 3d 规范场 $A$ 对偶化为周期标量，同时把周期标量 $A_4$ 对偶化为 3d 规范场。容易看出，电磁对偶作用在 A-型规范理论上给出以 $T^*[2]\C^*$ 为靶空间的 Rozansky-Witten 模型，即它把 A-型规范扇区映到 B-型物质扇区。类似地，它把 A-型物质扇区映到 B-型规范理论（规范群为 $U(1)$）。换言之，电磁对偶归结为粒子-涡旋对偶做两次。

% [Tutor Note: p.73 的记号 $T^*[2]\C^*$ 中，$[2]$ 表示把纤维坐标形式地赋予鬼数 2（对照 §6.1 对 $\C[2]$ 的说明：$\sigma$ 处在 2 次，故 $\C[2]$ 是纯偶的分次流形），不是复维数或实维数；因此 $T^*[2]\C^*$ 是分次辛流形，其上的 B-型 2d 理论就是 Rozansky-Witten 模型（见附录 B）。核对：本页 A-侧用 $T^*S^1$、B-侧用 $T^*[2]\C^*$，两种记号不要混用。]
